\documentclass[12pt]{article}
\usepackage[ngerman]{babel}\usepackage[babel=true]{microtype}
\usepackage[utf8]{inputenc}
\usepackage[T1]{fontenc}\usepackage{lmodern}\usepackage{color}
\usepackage[a4paper,includeheadfoot,
top=2cm, 
bottom=1cm, 
left=2cm, 
right=2cm]{geometry}
\usepackage{fancyhdr}\pagestyle{fancy}
\usepackage{placeins}
\usepackage{tabularx}
\usepackage{graphicx}
\usepackage{amsmath,amssymb,amsfonts,amsthm,mathtools}
\rhead{\vspace{-40pt}\includegraphics[scale=0.6]{logofsvwl.pdf}}
\cfoot{}
\lhead{}

 \def\chopline#1 \\{%für zwei Parameter \def\chopline#1;#2 \\{
  \def\name{#1}%\vorname{#2}
 }
 %\renewcommand{\familydefault}{\sfdefault}
 \newif\ifmore\moretrue
 
\newcolumntype{L}[1]{>{\raggedright\arraybackslash}p{#1}} % linksbündig mit Breitenangabe
\newcolumntype{C}[1]{>{\centering\arraybackslash}p{#1}} % zentriert mit Breitenangabe
\newcolumntype{R}[1]{>{\raggedleft\arraybackslash}p{#1}} % rechtsbündig mit Breitenangabe
\newcommand{\ltab}{\raggedright\arraybackslash} % Tabellenabschnitt linksbündig
\newcommand{\ctab}{\centering\arraybackslash} % Tabellenabschnitt zentriert
\newcommand{\rtab}{\raggedleft\arraybackslash} % Tabellenabschnitt rechtsbündig
\renewcommand{\labelitemii}{$\blacksquare$}
\renewcommand{\labelitemiii}{$\square$}
%
%
%
\begin{document} 
\newread\quelle
\openin\quelle=Teilnehmer.txt
\loop\read\quelle to \zeile\ifeof\quelle\global\morefalse
\else\expandafter\chopline\zeile\\
%%%%%%%%%%%%%%%%%%%%%%%%%%%%%%%%%%%%%%%%%%%%%%%%%%
%                  Seite 1                       %
%%%%%%%%%%%%%%%%%%%%%%%%%%%%%%%%%%%%%%%%%%%%%%%%%%
\begin{center}~\\[1cm]
{\fontsize{40}{0} \selectfont \textsf{\textbf{Bestätigung}}} \\[2cm]
   
{\fontsize{20}{0} \selectfont \textbf{\name}} \\[2cm]
   
{\fontsize{16}{0} \selectfont hat an der Lehrveranstaltung\\~\\der Fachschaft Volkswirtschaftslehre }\\[2cm]

   
{\fontsize{22}{0} \selectfont \textbf{\glqq Einführungskurs \LaTeX\grqq}}\\[2cm]
   
{\fontsize{16}{0} \selectfont in der Zeit vom 08.06 bis 10.06.2018 teilgenommen.} \\[2cm]
   
{\fontsize{16}{0} \selectfont Die Lehrveranstaltung umfasste insgesamt 20\\~\\ Unterrichtsstunden à 45 min} \\[3cm]
\begin{flushleft}
Bonn, den \today
\end{flushleft}~\\[0.5cm]
\FloatBarrier
\vspace{1,5cm}

\begin{tabular}{ccc} \cline{1-1} \cline{3-3}
Fachschaft  & ~~~~~~~~~~~~~~ & Dennis Kubitza\\
 Volkswirtschaftslehre & ~~~~~~~~~~~~~~ & (Dozent) \\
\end{tabular}

\end{center}
\newpage

%%%%%%%%%%%%%%%%%%%%%%%%%%%%%%%%%%%%%%%%%%%%%%%%%%
%                  Seite 2                       %
%%%%%%%%%%%%%%%%%%%%%%%%%%%%%%%%%%%%%%%%%%%%%%%%%%
~\\
\noindent
{\fontsize{20}{0} \selectfont \textbf{\underline{Inhalte der Veranstaltung:}}}\\
\begin{itemize}
\item {\large Grundlagen}
\begin{itemize}
\item Verwendung von Texmaker
\item Verwendung von Präambeln
\item Einbindung von Programmpaketen

\end{itemize}
\item {\large Textsatz}
	\begin{itemize}
	\item Textauszeichnung
	\item Einrichtung des Seitenformates
	\item Strukturierung von Dokumenten
		\begin{itemize}
		\item Erstellen eines Inhaltsverzeichnisses
		\item Veränderung der Seitennummerierung
		\item Verwendung der Kopf- und Fußzeile
		\item Verwendung von Gliederungsebenen 
		\item Erstellung von Aufzählungen
		\item Erstellung von Tabellen
		\item Einbindung von Grafiken
		\end{itemize}
	\end{itemize}
	
\item{\large Wissenschaftliches Schreiben}
\begin{itemize}
	\item {Erstellung von Theoremen und Definitionen}
	\item {Setzung mathematischer Formeln}
	\item {Verwendung interner Referenzen}
	\item {Verwendung von Bibliographien}
\end{itemize}
\end{itemize}

%%%%%%%%%%%%%%%%%%%%%%%%%%%%%%%%%%%%%%%%%%%%%%%%%%
%               Nächstes Dokument                %
%%%%%%%%%%%%%%%%%%%%%%%%%%%%%%%%%%%%%%%%%%%%%%%%%%
\newpage
\setcounter{page}{1}

%%%%%%%%%%%%%%%%%%%%%%%%%%%%%%%%%%%%%%%%%%%%%%%%%%
%               Ende des Dokuments               %
%%%%%%%%%%%%%%%%%%%%%%%%%%%%%%%%%%%%%%%%%%%%%%%%%%
% Eingabe darf keine Leerzeilen und keine letzte leere Zeile enthalten!
\fi\ifmore\repeat\closein\quelle
 
\end{document}